\chapter{Of Judicial and the Courts}
\label{art:judicial}

\articlenote{In which is set down the charge of Judicial: to rule where this Pact is silent, to judge where it is broken, and to keep, in one place, the record of how both are done.}

\section{The Charge and the Judge}

Judicial is charged with the judging of every offense under this Pact, the ruling of every matter this Pact leaves silent, and the keeping of the silo's Judicial archive, and holds its offices Up Top as Article 1 sets down.

\begin{clauses}
\item Judicial is led by a Judge, confirmed in office by the Mayor and, thereafter, by the Assembly under Article 20, and holding that office until death, resignation, or removal by the same two bodies acting together.
\item No citizen may hold the office of Judge who has not first served among the Officers of Judicial under Section 3, nor may a Judge also hold the office of Mayor, Sheriff, or Head of any other Department while serving.
\item The office of Judge is barred to a citizen afflicted with the Syndrome, and vacated by it, as Article 13, Section 6 provides; upon such a resignation the Mayor and the Assembly confirm a new Judge as this section provides.
\end{clauses}

\section{Rulings Where This Pact Is Silent}

Where this Pact is silent on a matter properly brought before it, the Judge shall issue a ruling in writing, stating the matter, the ruling, and the reasoning that led to it, and that ruling shall thereafter bind as though written into this Pact, exactly as Article 1 provides.

\begin{clauses}
\item A ruling under this section shall be kept in the Judicial archive and read into the record at the next sitting of the Assembly under Article 20, that no ruling binds the silo in secret.
\item A ruling under this section may be overturned only by a later ruling of the Judge on the same matter, or by amendment of this Pact directly under Article 23; it may not be overturned by the Mayor, the Sheriff, or any Head of Department.
\end{clauses}

\section{The Officers of Judicial}

Judicial shall keep such Officers as its charge requires, to investigate an offense, to gather what evidence the offense leaves behind, to search where a search is warranted, and to carry out a sentence once given, including the carrying of a citizen to cleaning under Article 18.

\begin{clauses}
\item An Officer of Judicial acts always at the direction of the Judge, or, in a matter touching the whole silo's safety, at the implicit command of the Mayor, and never on an Officer's own authority alone.
\item The Sheriff and the Officers of Judicial are not the same body. The Sheriff keeps the ordinary peace of the floors under Article 7; the Officers of Judicial act where Judicial itself has found cause to investigate, to search, or to carry out a sentence.
\item A search ordered by the Judge for a relic concealed under Article 16 may be made of any floor or room without further condition, save the station of the Sheriff, which may be searched only upon extreme cause found by the Judge, and only after notice of the search is given to the Sheriff beforehand.
\end{clauses}

\section{Judgment and Sentence}

Upon an offense investigated under Section 3, the Judge shall determine whether this Pact was broken and, if so, shall set the sentence this Pact provides, without further trial, and that sentence takes effect at once.

\begin{clauses}
\item The offender, or the offender's family where the offender cannot be heard, may make representations to the Judge after the sentence is set; the Judge shall hear them, but a representation made under this clause does not delay, lessen, or undo a sentence already taking effect.
\item A citizen accused of an offense not among the gravest named under Article 17 may request a hearing before sentence is set, and the Judge shall grant it; this right does not extend to an offense named among the gravest, for which Article 17 and this Article provide their own swiftness. In such a hearing, the accused citizen may speak, may call witnesses to speak to the accused's character or the facts of the offense, and may dispute the Officers' account of what occurred; the Judge shall hear all such evidence before determining guilt and setting sentence.
\item Ignorance of this Pact, where it excuses a first offense under Article 17, is weighed by the Judge before any sentence is set, not after.
\end{clauses}

\section{The Archive of Judicial}

Judicial shall keep, in its archive, a copy of this Pact, the rolls of citizenship under Article 2, every ruling issued under Section 2, every sentence given under Section 4, and such further record as this Pact elsewhere assigns to its keeping.

\begin{clauses}
\item No record kept in the Judicial archive may be altered once entered, save by a further entry noting the correction, the citizen who made it, and the reason given; the original entry shall not be removed.
\end{clauses}

\section{The Bringing of a Matter Before Judicial}

Every matter Judicial hears comes to it by one of three roads: a report of an offense, from the Sheriff under Article 7, from a Head of Department, or from any citizen; a petition, from any citizen of majority, upon a matter this Pact gives Judicial to decide; or a referral, from an officer this Pact directs to refer. This section sets down how each is entered.

\begin{clauses}
\item A report or petition is made to the Clerk of Judicial, in writing where the citizen can write and by the Clerk's hand where the citizen cannot, and is entered in the docket of Judicial with the day and hour of its entry, the name of the citizen bringing it, and the matter in a line. A citizen may not be refused the entry of a matter, though Judicial may afterward find it no matter for Judicial at all.
\item A report of an offense named among the gravest under Article 17 is carried to the Judge the same day; every other matter is taken in the docket's order, save that the Judge may take a matter out of its order where delay would work harm, and shall enter the reason.
\item No matter waits in the docket longer than one turn of the season from its entry without the Judge entering the cause of the delay; a citizen whose matter has waited longer may ask the Clerk the reason, and shall be told it.
\item A citizen who brings a matter may withdraw it before it is heard, and the withdrawal is entered; but a report of an offense, once entered, is Judicial's to pursue or to leave, and its withdrawal by the citizen who brought it does not end it.
\end{clauses}

\section{The Conduct of a Hearing}

A hearing under Section 4 is held in the hall of Judicial Up Top, or on such floor as the Judge names where the accused or the witnesses cannot climb to it, on a day the Clerk posts at the door of Judicial and at the landing of the accused's floor not fewer than ten days before.

\begin{clauses}
\item Present at a hearing are the Judge, the Clerk, the accused, the Officers who investigated, such witnesses as the accused or the Officers call, and the accused's household. Any other citizen may attend and may not be barred, save where the press of citizens endangers the hall, and save where the matter touches a relic or a thing this Pact keeps from ordinary record, in which case the Judge may clear the hall of all but those named.
\item A witness speaks upon the citizen's oath under Article 2, and is reminded of it before speaking; a witness found to have spoken falsely is answerable under Article 17, Section 5, as for a false bringing.
\item The Clerk keeps the record of every word spoken, so far as the Clerk's hand allows, and reads the record back before the Judge rules, that the accused may say whether it is true; the record and the accused's answer to it are entered in the archive together.
\item The Judge rules aloud and states the measure. The Clerk gives the accused a copy of the ruling in writing before the accused leaves the hall, and sends a copy to the Head of the accused's Department, and to Housing under Article 14 where the measure requires the accused's removal from a dwelling.
\item No hearing is held for an offense among the gravest; for such an offense Section 4 provides, and the Judge's ruling upon it is entered, read to the offender by the Clerk, and a copy given to the household. The representation Section 4 allows is heard in the same hall, on a day the Judge sets, and entered in the archive beside the ruling it does not alter.
\end{clauses}

\section{The Warrant and the Return of a Search}

An Officer of Judicial searches only upon a warrant of the Judge, save in the pursuit of an offender in the act, and a warrant names what it permits.

\begin{clauses}
\item A warrant is written and signed by the Judge, names the floor and room to be searched and the thing or citizen sought, and is shown to the citizen whose room is searched, or to the household, before the search begins; where none is present, it is left in the room.
\item Upon the search's end the Officer makes a return: what was searched, what was found, what was taken, and the citizen or Officer in whose keeping each thing taken now lies. The return is entered in the archive within one day, and the citizen searched may read it upon asking.
\item A thing taken under a warrant and found to be no relic, no restricted instrument, and no evidence is returned to the citizen within ten days, and the return is entered; a thing kept is classified under Article 16 or entered as evidence, and the citizen is told which.
\item A search made without a warrant, save in pursuit of an offender in the act, is a wrong under Article 17, Section 5, whatever it finds; but what it finds is not thereby unfound, and Judicial may act upon it.
\end{clauses}

\section{The Sittings of Judicial and the Clerk}

Judicial sits on posted days and keeps its own house by a Clerk, that no citizen need guess when a matter may be brought or by whom it is kept.

\begin{clauses}
\item The Judge sits to hear matters on no fewer than three days in every ten, on days posted at the door of Judicial; on other days the Judge rules upon petitions in writing and reads the docket.
\item The Clerk of Judicial keeps the docket, the record of hearings, and the Judicial archive under Section 5, and is raised by full shadowing under Article 13 from among the Officers or the archivists of Judicial; the Clerk does not rule.
\item Where the Judge is absent or unable to sit, the senior Officer of Judicial by first confirmation keeps the docket and may order a search under Section 3, but may neither rule nor sentence; a matter that cannot wait is held, and the Judge's absence beyond one turn of the season is reported to the Mayor and to the Assembly as a vacancy under Section 1.
\end{clauses}

\section{The Petition for a Ruling Where This Pact Is Silent}

A ruling under Section 2 is not given unasked; it is given upon a matter properly brought, and this section sets down how such a matter is brought and argued.

\begin{clauses}
\item Any citizen of majority, any office, or the Judge upon the Judge's own reading of the docket, may bring a matter under Section 2 by a petition entered with the Clerk under Section 6, stating the matter, the Article or Articles the petitioner believes nearest to it, and what the petitioner believes this Pact would provide if it spoke.
\item The Clerk posts the petition at the door of Judicial for one cycle, and any citizen may enter with the Clerk, within that cycle, a writing in answer to it, stating what that citizen believes this Pact would provide; a citizen who enters such a writing is heard at the ruling if the citizen asks to be.
\item The Judge hears the petitioner and every citizen entered under this section on a posted day, in the hall of Judicial, and may put questions to each; the Judge then rules in writing under Section 2, and the writing states which of the readings offered the Judge took, and why, or that the Judge took none of them.
\item A matter brought under this section while an offense turning upon it stands in the docket is ruled before the offense is heard, and the ruling governs the offense; but no ruling under Section 2 reaches back to make an offense of a thing done before the ruling was entered.
\item The rulings under Section 2 are kept in their own roll in the archive, in the order of their entry, and the Clerk keeps beside it a table of the Articles each ruling touches, that a citizen may find what this Pact has been held to provide upon any matter; the table is open to any citizen at the door of Judicial.
\end{clauses}

\section{Evidence and Things Seized}

What an Officer takes under a warrant, and what a citizen brings to Judicial in proof, is kept in a manner that lets it be found again and shows it has not been changed.

\begin{clauses}
\item Every thing taken under a warrant under Section 8, or surrendered under Article 16, or brought to Judicial by a citizen as proof of a matter, is entered by the Clerk in the roll of evidence with the day, the matter it belongs to, the citizen from whom or by whom it came, and a description sufficient to know it again; the Clerk marks the thing with the number of its entry.
\item A thing so entered is kept in the store of Judicial under lock, and is taken out only by an Officer or the Clerk, who enters the taking-out, the purpose, and the return; a thing that is a relic or a restricted instrument passes from the store to Information Technology under Article 8 for examination, and the passing is entered, and Information Technology's report under that Article is entered beside it.
\item A thing entered in the roll of evidence is shown to the citizen it was taken from, or to the accused in whose matter it is proof, at the hearing under Section 7, or, where there is no hearing, upon the citizen's asking at the door of Judicial; the citizen may say what the thing is and how it came to be where it was found, and the Clerk enters what the citizen says.
\item A thing that is neither a relic, a restricted instrument, nor proof of an offense found is returned as Section 8 provides; a thing that is proof is kept until the matter is ruled and for one year after, and then returned, or, where it is the silo's or cannot be returned, given to Mechanical for reclamation under Article 9 and the giving entered.
\item A thing that is a relic or a restricted instrument is never returned save under a sanction given under Article 16, and is destroyed, sealed, or kept as that Article and Article 8 provide, the disposal entered in the roll of evidence and in the relic roll both.
\end{clauses}

\section{The Entry and Carrying Out of a Sentence}

A sentence set under Section 4 is a written thing before it is anything else, and it is carried out as it is written and not otherwise.

\begin{clauses}
\item The Clerk enters every sentence in the roll of sentences with the citizen's name, the offense, the Article and Section under which it is answered, the measure, and, for a term in the mines, the span; the Judge signs the entry, and no sentence is carried out before it is signed.
\item A grace under Article 17 is entered and a copy given to the citizen; no further act is required, and the citizen returns to trade and dwelling the same day.
\item A term in the mines is carried out by two Officers of Judicial, who conduct the citizen by the stair to the Down Deep within one cycle of the sentence, deliver the citizen to the Head of Mines with a copy of the entry, and take the Head's receipt for the citizen upon the copy; the receipt is entered in the archive. The Head of Mines reports the term's ending to Judicial under Article 10, and the Clerk enters it.
\item A citizen returned from the mines presents before the Clerk within one cycle of the term's ending, and Judicial assigns the citizen's trade and floor under Article 17 and reports the assignment to the citizen's new Head of Department and to Housing under Article 14; a citizen so returned holds every standing of a citizen, and no office may refuse the citizen a trade, a dwelling, or the ration on account of the term served.
\item A sentence of cleaning is carried out as Article 18 provides, and the Officer or Sheriff who conducts the citizen to the airlock returns to the Clerk with the chamber's record, which is entered beside the sentence.
\item Restitution ordered under Article 17, Section 5 is entered with the sentence, and the Chit-Keepers under Article 12 are given a copy where it is in chits, and the Head of the offender's Department where it is in labor; the Clerk enters its completion when the wronged citizen or the Chit-Keepers report it.
\end{clauses}

\section{The Officers and Their Discipline}

An Officer of Judicial holds a citizen's liberty in hand more often than any other citizen of the silo, and answers for it more closely.

\begin{clauses}
\item An Officer of Judicial is entered by full shadowing under Article 13 and the vetting that Article provides for a critical role, and is confirmed by the Judge; an Officer serves at the Judge's direction and may be removed by the Judge for cause entered in the archive.
\item An Officer acts only under a warrant, a sentence, or the Judge's written direction, save in the pursuit of an offender in the act, and an Officer who acts otherwise is answerable under Article 17 as any citizen, and the act is a wrong under Section 5 of that Article, whatever it found.
\item A citizen who believes an Officer to have exceeded a warrant, to have taken what the warrant did not name, or to have done the citizen a wrong in the course of a search, an arrest, or the carrying out of a sentence, may bring a complaint against the Officer under Section 6, and the Judge shall hear it as any other matter; the Officer complained of does not investigate it, and the Judge names another to do so.
\item A complaint against an Officer found true is entered against the Officer's name, and the Judge sets the measure under Article 17; a complaint found false is entered against the complainant as a false bringing. Two complaints found true against one Officer remove the Officer from the Officers, whatever measure was set for each.
\item The Officers keep no record apart from the archive: every warrant, return, arrest, and delivery an Officer makes is entered with the Clerk, and an Officer who keeps a record of the Officer's own, apart from the Clerk's, is removed.
\end{clauses}
